% !PRINTED 99
\weberpage{99}
%<<CONT

folgt. In dem Product aller dieser Ausdrücke kommt jeder Factor des Productes (1) zweimal mit entgegengesetztem Zeichen vor, so dass wir, da die Anzahl der Factoren von (1) gleich $\frac{n(n-1)}{2}$ ist, die Gleichung erhalten
\begin{equation}
 f'(\alpha_1)f'(\alpha_2)\cdots f'(\alpha_n)=(-1)^{\frac{n(n-1)}2}[\alpha_1,\alpha_2,\ldots,\alpha_n]^2.\tag{5}
\end{equation}
Wir können unser Product (1) aber auch mit Hülfe der Relationen (4) folgendermaassen darstellen:
\begin{align}
 [\alpha_1,\alpha_2,\ldots,\alpha_n]
 &= f'(\alpha_1)[\alpha_2,\alpha_3,\ldots,\alpha_n]\notag\\
 &=-f'(\alpha_2)[\alpha_1,\alpha_3,\ldots,\alpha_n]\tag{6}\\
 &=+f'(\alpha_3)[\alpha_1,\alpha_2,\ldots,\alpha_n]\notag\\% [sic] Weber druckt hier die vollstaendige Folge alpha_1,alpha_2,...,alpha_n, obwohl der Faktor f'(alpha_3) die Klammer ohne alpha_3 erfordert und der unmittelbar folgende Satz ausdruecklich sagt, jeder rechte Klammerausdruck enthalte ein Element weniger. Quellentreuer mathematischer Druckfehler; nicht stillschweigend zu [alpha_1,alpha_2,alpha_4,...,alpha_n] normalisiert.
 &=\cdots=(-1)^{n-1}f'(\alpha_n)[\alpha_1,\alpha_2,\ldots,\alpha_{n-1}],\notag
\end{align}
wo auf der rechten Seite jeder der Klammerausdrücke ein Element weniger enthält als der Klammerausdruck auf der linken Seite.

Es soll also jetzt die Aufgabe gestellt sein, eine ganze rationale Function $(n-1)$ten Grades $\varphi(x)$ zu bestimmen, die für die Argumentwerthe $x=\alpha_1,\alpha_2,\ldots,\alpha_n$ der Reihe nach die vorgeschriebenen Werthe $\varphi(\alpha_1),\varphi(\alpha_2),\ldots,\varphi(\alpha_n)$ annimmt.

Setzt man
\begin{equation}
 \varphi(x)=a_1x^{n-1}+a_2x^{n-2}+\cdots+a_{n-1}x+a_n,\tag{7}
\end{equation}
so sind die $n$ unbekannten Coëfficienten $a_1,a_2,\ldots,a_n$ aus den linearen Gleichungen
\begin{align}
 \varphi(\alpha_1)&=a_1\alpha_1^{n-1}+a_2\alpha_1^{n-2}+\cdots+a_{n-1}\alpha_1+a_n,\notag\\
 \varphi(\alpha_2)&=a_1\alpha_2^{n-1}+a_2\alpha_2^{n-2}+\cdots+a_{n-1}\alpha_2+a_n,\tag{8}\\
 &\vdots\notag\\
 \varphi(\alpha_n)&=a_1\alpha_n^{n-1}+a_2\alpha_n^{n-2}+\cdots+a_{n-1}\alpha_n+a_n\notag
\end{align}
zu bestimmen. Die Determinante dieses Systems ist aber genau die Grösse (2)
\[
[\alpha_1,\alpha_2,\ldots,\alpha_n],
\]
und die Aufgabe ist also nach §. 26 immer lösbar, wenn diese Grösse von Null verschieden ist, d. h. wenn von den $n$ Grössen $\alpha_1,\alpha_2,\ldots,\alpha_n$ keine zwei einander gleich sind.

Statt aber die Gleichungen (8) nach §. 26 aufzulösen und die gefundenen Ausdrücke in (7) einzusetzen, ist es vorzuziehen, nach §. 24, II. die homogenen Grössen
\[
1,a_1,a_2,\ldots,a_n
\]
%CONT>>
